\textbf{Hysterezní smyčka} vzniká při cyklické změně pole; magnetický stav závisí i na předchozím magnetování.
\begin{itemize}
\item Z nasycení při poklesu na $H=0$ zbývá \textbf{remanence} $B_r$ nebo $M_r$.
\item Opačné \textbf{koercitivní pole} $-H_c$ zruší příslušnou indukci či magnetizaci; silnější opačné pole vytvoří záporné nasycení.
\item Zpětná větev přes zápornou remanenci a kladnou koercivitu uzavírá smyčku.
\end{itemize}
\textbf{Křivka prvotní magnetizace} naproti tomu není uzavřená: vede z odmagnetovaného počátku k nasycení při postupném zvyšování pole. Hysterezní smyčka popisuje celý cyklus přemagnetování.
